\section*{Earth's Green Mantle / 大地的绿色披覆}
\begin{flushright}
Rachel Carson, \textit{Silent Spring}
\end{flushright}

\noindent\textbf{English.} \ldots{} the booming ``weed killer'' business of the present day, in which soaring sales and expanding uses mark the production of plant-killing chemicals.
\par
\noindent\textbf{中文。} ……当今蓬勃发展的“除草剂”产业：销量节节攀升，用途不断扩大，推动着各种杀灭植物的化学制剂大量生产。
\par\medskip

\noindent\textbf{English.} One of the most tragic examples of our unthinking bludgeoning of the landscape is to be seen in the sagebrush lands of the West, where a vast campaign is on to destroy the sage and to substitute grasslands. If ever an enterprise needed to be illuminated with a sense of the history and meaning of the landscape, it is this. For here the natural landscape is eloquent of the interplay of forces that have created it. It is spread before us like the pages of an open book in which we can read why the land is what it is, and why we should preserve its integrity. But the pages lie unread.
\par
\noindent\textbf{中文。} 我们不假思索地重击大地景观，其中最悲惨的例子之一，就出现在西部的蒿灌丛地带：一场大规模运动正在进行，要消灭蒿灌木，代之以草地。倘若有哪项事业尤其需要以对景观历史与意义的理解来照亮，那正是这一项。因为这里的自然景观清楚讲述了塑造它的诸力如何相互作用。它像一本摊开的书展现在我们面前，使我们能够读出土地为何如此，也读出为什么应当保存它的完整性。可是这些书页无人阅读。
\par\medskip

\noindent\textbf{English.} The land of the sage is the land of the high western plains and the lower slopes of the mountains that rise above them, a land born of the great uplift of the Rocky Mountain system many millions of years ago. It is a place of harsh extremes of climate: of long winters when blizzards drive down from the mountains and snow lies deep on the plains, of summers whose heat is relieved by only scanty rains, with drought biting deep into the soil, and drying winds stealing moisture from leaf and stem.
\par
\noindent\textbf{中文。} 蒿的土地，就是西部高原和从其上拔起的山脉低坡；这片土地诞生于数百万年前落基山系的巨大隆升。这里气候严酷而极端：漫长的冬季，暴风雪从山中压下，平原积雪深厚；夏季炎热，只有稀少雨水稍作缓解，干旱深深咬入土壤，干风从叶片和茎秆中偷走水分。
\par\medskip

\noindent\textbf{English.} As the landscape evolved, there must have been a long period of trial and error in which plants attempted the colonization of this high and windswept land. One after another must have failed. At last one group of plants evolved which combined all the qualities needed to survive. The sage---low-growing and shrubby---could hold its place on the mountain slopes and on the plains, and within its small gray leaves it could hold moisture enough to defy the thieving winds. It was no accident, but rather the result of long ages of experimentation by nature, that the great plains of the West became the land of the sage.
\par
\noindent\textbf{中文。} 在景观演化的过程中，植物试图定居这片高寒、多风的土地，必定经历过漫长的试错。一种又一种植物想必都失败了。最后，一类植物演化出来，兼具生存所需的一切品质。蒿灌木低矮而丛生，能够在山坡和平原上站稳脚跟，并能在小小的灰色叶片里保存足够的水分，抵抗偷水的风。西部大平原成为蒿的土地，并非偶然，而是自然经历漫长年代实验的结果。
\par\medskip

\noindent\textbf{English.} Along with the plants, animal life, too, was evolving in harmony with the searching requirements of the land. In time there were two as perfectly adjusted to their habitat as the sage. One was a mammal, the fleet and graceful pronghorn antelope. The other was a bird, the sage grouse---the ``cock of the plains'' of Lewis and Clark.
\par
\noindent\textbf{中文。} 与植物一道，动物生命也在演化，并与这片土地苛刻的要求相协调。久而久之，有两种动物像蒿一样完美地适应了它们的栖息地。一种是哺乳动物，迅捷而优雅的叉角羚；另一种是鸟类，蒿原松鸡，也就是刘易斯和克拉克所称的“平原雄鸡”。
\par\medskip

\noindent\textbf{English.} The sage and the grouse seem made for each other. The original range of the bird coincided with the range of the sage, and as the sagelands have been reduced, so the populations of grouse have dwindled. The sage is all things to these birds of the plains. The low sage of the foothill ranges shelters their nests and their young; the denser growths are loafing and roosting areas; at all times the sage provides the staple food of the grouse. Yet it is a two-way relationship. The spectacular courtship displays of the cocks help loosen the soil beneath and around the sage, aiding invasion by grasses which grow in the shelter of sagebrush.
\par
\noindent\textbf{中文。} 蒿和松鸡仿佛天生相属。蒿原松鸡最初的分布范围与蒿灌丛的范围相吻合；随着蒿地缩减，松鸡数量也随之下降。对这些平原之鸟来说，蒿就是一切。山麓地带低矮的蒿庇护它们的巢和幼雏；更浓密的蒿丛是它们闲栖和夜宿之所；无论何时，蒿都是松鸡的主食。然而，这也是一种双向关系。雄鸟壮观的求偶炫耀，会松动蒿下及其周围的土壤，帮助草类侵入并在蒿灌木庇护下生长。
\par\medskip

\noindent\textbf{English.} The antelope, too, have adjusted their lives to the sage. They are primarily animals of the plains, and in winter when the first snows come those that have summered in the mountains move down to the lower elevations. There the sage provides the food that tides them over the winter. Where all other plants have shed their leaves, the sage remains evergreen, the gray-green leaves---bitter, aromatic, rich in proteins, fats, and needed minerals---clinging to the stems of the dense and shrubby plants. Though the snows pile up, the tops of the sage remain exposed, or can be reached by the sharp, pawing hoofs of the antelope. Then grouse feed on them too, finding them on bare and windswept ledges or following the antelope to feed where they have scratched away the snow.
\par
\noindent\textbf{中文。} 叉角羚也把自己的生活调适到蒿的节律之中。它们主要是平原动物；冬天初雪降临时，在山地度过夏季的羚群便下移到较低海拔。那里，蒿提供食物，帮助它们熬过冬季。其他植物都已落叶，蒿却依然常绿，灰绿色的叶片苦涩、芳香，富含蛋白质、脂肪和必需矿物质，紧附在浓密丛生的茎上。即使积雪堆深，蒿梢仍会露出，或可由叉角羚锋利、刨挖的蹄子触及。于是松鸡也以其为食，在裸露而风扫的岩脊上寻找，或跟随叉角羚，在它们刨开雪的地方觅食。
\par\medskip

\noindent\textbf{English.} And other life looks to the sage. Mule deer often feed on it. Sage may mean survival for winter-grazing livestock. Sheep graze many winter ranges where the big sagebrush forms almost pure stands. For half the year it is their principal forage, a plant of higher energy value than even alfalfa hay.
\par
\noindent\textbf{中文。} 还有其他生命仰赖蒿。骡鹿常以它为食。对于冬季放牧的牲畜来说，蒿可能意味着生存。许多冬季牧场上，大蒿几乎形成纯林，羊群就在其中采食。一年中有半年，蒿是它们的主要饲草，其能量价值甚至高于苜蓿干草。
\par\medskip

\noindent\textbf{English.} The bitter upland plains, the purple wastes of sage, the wild, swift antelope, and the grouse are then a natural system in perfect balance. Are? The verb must be changed---at least in those already vast and growing areas where man is attempting to improve on nature's way. In the name of progress the land management agencies have set about to satisfy the insatiable demands of the cattlemen for more grazing land. By this they mean grassland---grass without sage. So in a land which nature found suited to grass growing mixed with and under the shelter of sage, it is now proposed to eliminate the sage and create unbroken grassland. Few seem to have asked whether grasslands are a stable and desirable goal in this region. Certainly nature's own answer was otherwise. The annual precipitation in this land where the rains seldom fall is not enough to support good sod-forming grass; it favors rather the perennial bunch grass that grows in the shelter of the sage.
\par
\noindent\textbf{中文。} 苦寒的高地平原、紫色的蒿海、野性而迅疾的叉角羚以及松鸡，原本构成一个完美平衡的自然系统。原本？这个动词必须改掉，至少在那些已经广阔且仍在扩大的地区是如此，因为人在那里试图改进自然之道。以进步之名，土地管理机构开始满足牧牛人对更多牧场无止境的要求。他们所谓的牧场，指的是草地，只有草、没有蒿。于是，在自然认为适宜草与蒿混生、并在蒿的庇护下生长的土地上，如今有人提议消灭蒿，创造连绵不断的草地。似乎很少有人问过，在这一地区，草地是否真是稳定而可取的目标。自然自身的回答显然并非如此。在这片少雨之地，年降水量不足以支撑结成良好草皮的禾草；它更有利于在蒿庇护下生长的多年生丛生禾草。
\par\medskip

\noindent\textbf{English.} Yet the program of sage eradication has been under way for a number of years. Several government agencies are active in it; industry has joined with enthusiasm to promote and encourage an enterprise which creates expanded markets not only for grass seed but for a large assortment of machines for cutting and plowing and seeding. The newest addition to the weapons is the use of chemical sprays. Now millions of acres of sagebrush lands are sprayed each year.
\par
\noindent\textbf{中文。} 然而，根除蒿的计划已经推行多年。若干政府机构积极参与；工业界也热情加入，推动和鼓励这项事业，因为它不仅为草籽，也为各种割灌、犁耕和播种机械开辟了更大的市场。最新加入武器库的是化学喷洒。如今，每年都有数百万英亩蒿灌丛地被喷洒药剂。
\par\medskip

\noindent\textbf{English.} What are the results? The eventual effects of eliminating sage and seeding with grass are largely conjectural. Men of long experience with the ways of the land say that in this country there is better growth of grass between and under the sage than can possibly be had in pure stands, once the moisture-holding sage is gone.
\par
\noindent\textbf{中文。} 结果如何？消灭蒿并播种牧草的最终影响，在很大程度上仍属推测。那些长期熟悉这片土地规律的人说，在这个地区，草在蒿丛之间和蒿下生长得更好；一旦能够保持水分的蒿消失，纯草地绝不可能达到同样的长势。
\par\medskip

\noindent\textbf{English.} But even if the program succeeds in its immediate objective, it is clear that the whole closely knit fabric of life has been ripped apart. The antelope and the grouse will disappear along with the sage. The deer will suffer, too, and the land will be poorer for the destruction of the wild things that belong to it. Even the livestock which are the intended beneficiaries will suffer; no amount of lush green grass in summer can help the sheep starving in the winter storms for lack of the sage and bitterbrush and other wild vegetation of the plains.
\par
\noindent\textbf{中文。} 但是，即使这项计划实现了眼前目标，也很清楚，整个紧密编织的生命之网已经被撕裂。叉角羚和松鸡将随蒿一起消失。鹿也会受害；那些原本属于这片土地的野生生命被毁灭，土地也将变得更贫乏。即便是本应受益的牲畜也会遭殃；夏季再茂盛的青草，也无法帮助羊群在冬季风雪中挨饿，因为它们失去了蒿、苦灌木以及平原上的其他野生植被。
\par\medskip

\noindent\textbf{English.} These are the first and obvious effects. The second is of a kind that is always associated with the shotgun approach to nature: the spraying also eliminates a great many plants that were not its intended target. Justice William O. Douglas, in his recent book \textit{My Wilderness: East to Katahdin}, has told of an appalling example of ecological destruction wrought by the United States Forest Service in the Bridger National Forest in Wyoming. Some 10,000 acres of sagelands were sprayed by the Service, yielding to pressure of cattlemen for more grasslands. The sage was killed, as intended. But so was the green, life-giving ribbon of willows that traced its way across these plains, following the meandering streams. Moose had lived in these willow thickets, for willow is to the moose what sage is to the antelope. Beavers had lived there, too, feeding on the willows, felling them and making a strong dam across the tiny stream. Through the labor of the beavers, a lake backed up. Trout in the mountain streams seldom were more than six inches long; in the lake they thrived so prodigiously that many grew to five pounds. Waterfowl were attracted to the lake, also. Merely because of the presence of the willows and the beavers that depended on them, the region was an attractive recreational area with excellent fishing and hunting.
\par
\noindent\textbf{中文。} 这些是最初且显而易见的影响。第二类影响，则总是与对自然采取霰弹枪式的粗暴办法相伴而来：喷洒还会消灭许多并非目标的植物。威廉・O・道格拉斯大法官在近作《我的荒野：从东部到卡塔丁》中，讲述了美国林务局在怀俄明州布里杰国家森林造成生态破坏的一个惊人例子。林务局屈从牧牛人要求更多草地的压力，喷洒了约一万英亩蒿地。蒿按计划被杀死了。但那条沿蜿蜒溪流横贯平原、带来生命的绿色柳带也一同死去。驼鹿曾生活在这些柳丛中，因为柳树之于驼鹿，正如蒿之于叉角羚。河狸也曾在那里生活，以柳为食，伐倒柳树，在小溪上筑起坚固堤坝。由于河狸的劳动，水积成了湖。山溪中的鳟鱼很少超过六英寸；在湖中，它们却繁盛异常，许多长到五磅。水禽也被湖吸引而来。仅仅因为有了柳树和依赖柳树的河狸，这一地区便成为有优良垂钓与狩猎条件、富有吸引力的休闲地。
\par\medskip

\noindent\textbf{English.} But with the ``improvement'' instituted by the Forest Service, the willows went the way of the sagebrush, killed by the same impartial spray. When Justice Douglas visited the area in 1959, the year of the spraying, he was shocked to see the shriveled and dying willows---the ``vast, incredible damage.'' What would become of the moose? Of the beavers and the little world they had constructed? A year later he returned to read the answers in the devastated landscape. The moose were gone and so were the beavers. Their principal dam had gone out for want of attention by its skilled architects, and the lake had drained away. None of the large trout were left. None could live in the tiny creek that remained, threading its way through a bare, hot land where no shade remained. The living world was shattered.
\par
\noindent\textbf{中文。} 然而，随着林务局推行的“改良”，柳树也步了蒿灌木的后尘，被同一种不分青红皂白的喷剂杀死。道格拉斯大法官在1959年，也就是喷洒当年访问该地时，看到枯缩将死的柳树，看到那种“巨大而难以置信的损害”，深感震惊。驼鹿会怎样？河狸和它们构筑的那个小世界又会怎样？一年后，他重返此地，在被毁坏的景观中读到了答案。驼鹿不见了，河狸也不见了。它们的主坝因缺少熟练建筑师的维护而垮掉，湖水也流干。大型鳟鱼一条也没有留下。残存的小溪穿过光秃、炎热、毫无遮阴的土地，没有一条大鳟鱼能在那里活下去。这个有生命的世界被粉碎了。
\par\medskip

\noindent\textbf{English.} Besides the more than four million acres of rangelands sprayed each year, tremendous areas of other types of land are also potential or actual recipients of chemical treatments for weed control. For example, an area larger than all of New England---some 50 million acres---is under management by utility corporations and much of it is routinely treated for ``brush control.'' In the Southwest an estimated 75 million acres of mesquite lands require management by some means, and chemical spraying is the method most actively pushed. An unknown but very large acreage of timber-producing lands is now aerially sprayed in order to ``weed out'' the hardwoods from the more spray-resistant conifers. Treatment of agricultural lands with herbicides doubled in the decade following 1949, totaling 53 million acres in 1959. And the combined acreage of private lawns, parks, and golf courses now being treated must reach an astronomical figure.
\par
\noindent\textbf{中文。} 除了每年被喷洒的四百多万英亩牧地之外，其他类型的巨大面积土地，也正在或可能正在接受化学除草处理。比如，由公用事业公司管理的土地面积超过整个新英格兰，约五千万英亩，其中很大一部分被例行用于“灌丛控制”。在西南部，估计有七千五百万英亩牧豆树地需要以某种方式管理，而化学喷洒是被最积极推行的方法。另有面积不明但极为广阔的产材林地，如今正通过空中喷洒来从较耐喷剂的针叶树中“剔除”阔叶树。1949年以后的十年间，农地使用除草剂处理的面积翻了一番，到1959年总计达五千三百万英亩。至于如今接受处理的私人草坪、公园和高尔夫球场，合计面积必定达到天文数字。
\par\medskip

\noindent\textbf{English.} The chemical weed killers are a bright new toy. They work in a spectacular way; they give a giddy sense of power over nature to those who wield them, and as for the long-range and less obvious effects---these are easily brushed aside as the baseless imaginings of pessimists. The ``agricultural engineers'' speak blithely of ``chemical plowing'' in a world that is urged to beat its plowshares into spray guns. The town fathers of a thousand communities lend willing ears to the chemical salesman and the eager contractors who will rid the roadsides of ``brush''---for a price. It is cheaper than mowing, is the cry. So, perhaps, it appears in the neat rows of figures in the official books; but were the true costs entered, the costs not only in dollars but in the many equally valid debits we shall presently consider, the wholesale broadcasting of chemicals would be seen to be more costly in dollars as well as infinitely damaging to the long-range health of the landscape and to all the varied interests that depend on it.
\par
\noindent\textbf{中文。} 化学除草剂是一件闪亮的新玩具。它们起效的方式令人瞩目；使用者会因此产生一种对自然拥有权力的眩晕感。至于长期而不那么明显的后果，则很容易被斥为悲观者毫无根据的想象。“农业工程师”轻快地谈论“化学耕作”，仿佛世界正被敦促把犁铧改铸成喷枪。成千上万社区的镇政要人乐于听信化学品推销员和急切承包商的话；这些人愿意替道路两旁清除“灌丛”，只要付钱即可。呼声说，这比割草便宜。也许在官方账册整齐排列的数字中，看起来确是如此；可是如果把真实成本列入账内，不仅是美元成本，也包括我们马上要考量的那些同样有效的借项，那么，大规模撒布化学品就会显得在金钱上更昂贵，并且对景观的长期健康、对依赖景观的各种利益造成无穷损害。
\par\medskip

\noindent\textbf{English.} Take, for instance, that commodity prized by every chamber of commerce throughout the land---the good will of vacationing tourists. There is a steadily growing chorus of outraged protest about the disfigurement of once beautiful roadsides by chemical sprays, which substitute a sere expanse of brown, withered vegetation for the beauty of fern and wildflower, of native shrubs adorned with blossom or berry. ``We are making a dirty, brown, dying-looking mess along the sides of our roads,'' a New England woman wrote angrily to her newspaper. ``This is not what the tourists expect, with all the money we are spending advertising the beautiful scenery.''
\par
\noindent\textbf{中文。} 例如，全国每个商会都珍视的一项资产：度假游客的好感。对化学喷洒毁坏昔日美丽路边景观的愤怒抗议，正日益汇成合唱；喷洒以一片焦枯、褐色、萎败的植被，取代了蕨类和野花之美，取代了开花或结莓的本地灌木之美。一位新英格兰妇女愤怒地写信给报纸说：“我们正在把道路两旁弄成肮脏、褐色、垂死模样的一团糟。我们花了那么多钱宣传美丽风景，游客期待看到的可不是这个。”
\par\medskip

\noindent\textbf{English.} In the summer of 1960 conservationists from many states converged on a peaceful Maine island to witness its presentation to the National Audubon Society by its owner, Millicent Todd Bingham. The focus that day was on the preservation of the natural landscape and of the intricate web of life whose interwoven strands lead from microbes to man. But in the background of all the conversations among the visitors to the island was indignation at the despoiling of the roads they had traveled. Once it had been a joy to follow those roads through the evergreen forests, roads lined with bayberry and sweet fern, alder and huckleberry. Now all was brown desolation. One of the conservationists wrote of that August pilgrimage to a Maine island: ``I returned \ldots{} angry at the desecration of the Maine roadsides. Where, in previous years, the highways were bordered with wildflowers and attractive shrubs, there were only the scars of dead vegetation for mile after mile \ldots{}. As an economic proposition, can Maine afford the loss of tourist goodwill that such sights induce?''
\par
\noindent\textbf{中文。} 1960年夏，来自许多州的自然保护人士聚集到缅因州一座宁静小岛，见证岛主米利森特・托德・宾厄姆将其赠予全国奥杜邦学会。那一天的焦点，是保存自然景观，保存那张从微生物一直连到人类、千丝万缕相互交织的生命之网。然而，访客们在岛上所有交谈的背景里，都有一种对来路遭到破坏的愤慨。过去，沿着那些道路穿过常绿森林曾是一种喜悦，道旁生长着杨梅属灌木、香蕨木、桤木和越橘。如今，一切都是褐色荒凉。一位保护人士写到那年八月前往缅因小岛的旅程：“我回来时……为缅因路边遭到亵渎而愤怒。往年，公路边满是野花和迷人的灌木；如今一英里又一英里，只剩死去植被留下的伤疤……从经济角度看，缅因能承受这种景象导致的游客好感流失吗？”
\par\medskip

\noindent\textbf{English.} Maine roadsides are merely one example, though a particularly sad one for those of us who have a deep love for the beauty of that state, of the senseless destruction that is going on in the name of roadside brush control throughout the nation.
\par
\noindent\textbf{中文。} 缅因的路边不过是一个例子；对于我们这些深深爱着该州之美的人来说，这是一个尤其令人悲伤的例子。全国各地以道路灌丛控制之名进行的无谓破坏，正在不断发生。
\par\medskip

\noindent\textbf{English.} Botanists at the Connecticut Arboretum declare that the elimination of beautiful native shrubs and wildflowers has reached the proportions of a ``roadside crisis.'' Azaleas, mountain laurel, blueberries, huckleberries, viburnums, dogwood, bayberry, sweet fern, low shadbush, winterberry, chokecherry, and wild plum are dying before the chemical barrage. So are the daisies, black-eyed Susans, Queen Anne's lace, goldenrods, and fall asters which lend grace and beauty to the landscape.
\par
\noindent\textbf{中文。} 康涅狄格植物园的植物学家宣称，消灭美丽本地灌木和野花的做法，已经发展到“路边危机”的程度。杜鹃、山月桂、蓝莓、越橘、荚蒾、山茱萸、杨梅属灌木、香蕨木、矮唐棣、冬青、稠李和野李，都在化学弹幕前死去。雏菊、黑心金光菊、野胡萝卜、秋麒麟草和秋紫苑也同样如此；它们本来为景观增添优雅与美丽。
\par\medskip

\noindent\textbf{English.} The spraying is not only improperly planned but studded with abuses such as these. In a southern New England town one contractor finished his work with some chemical remaining in his tank. He discharged this along woodland roadsides where no spraying had been authorized. As a result the community lost the blue and golden beauty of its autumn roads, where asters and goldenrod would have made a display worth traveling far to see. In another New England community a contractor changed the state specifications for town spraying without the knowledge of the highway department and sprayed roadside vegetation to a height of eight feet instead of the specified maximum of four feet, leaving a broad, disfiguring, brown swath. In a Massachusetts community the town officials purchased a weed killer from a zealous chemical salesman, unaware that it contained arsenic. One result of the subsequent roadside spraying was the death of a dozen cows from arsenic poisoning.
\par
\noindent\textbf{中文。} 喷洒不仅规划不当，还充斥着这类滥用。在新英格兰南部一座城镇，一名承包商完工后罐中还剩一些药剂，便把它沿林地路边排放，而那里并未获准喷洒。结果，这个社区失去了秋日道路的蓝金之美；原本紫苑和秋麒麟草会在那里形成一场值得远道而来的展示。在另一个新英格兰社区，一名承包商在公路部门不知情的情况下改变了州里关于城镇喷洒的规格，把路边植被喷到八英尺高，而规定的最高限度是四英尺，留下了一条宽阔、丑陋的褐色带状伤痕。在马萨诸塞州一个社区，镇官员从一名热心的化学品推销员那里购买除草剂，却不知道其中含砷。随后进行的路边喷洒，造成十几头牛死于砷中毒。
\par\medskip

\noindent\textbf{English.} Trees within the Connecticut Arboretum Natural Area were seriously injured when the town of Waterford sprayed the roadsides with chemical weed killers in 1957. Even large trees not directly sprayed were affected. The leaves of the oaks began to curl and turn brown, although it was the season for spring growth. Then new shoots began to be put forth and grew with abnormal rapidity, giving a weeping appearance to the trees. Two seasons later, large branches on these trees had died, others were without leaves, and the deformed, weeping effect of whole trees persisted.
\par
\noindent\textbf{中文。} 1957年，沃特福德镇在路边喷洒化学除草剂时，康涅狄格植物园自然区内的树木受到严重伤害。即使没有被直接喷到的大树也受到影响。橡树叶片开始卷曲、变褐，尽管那正是春季生长时节。随后，新梢开始抽出，并以异常速度生长，使树呈现垂泣般的外观。两个生长季之后，这些树上的大枝已经死亡，另一些枝条没有叶片，整株树畸形下垂的状态仍然存在。
\par\medskip

\noindent\textbf{English.} I know well a stretch of road where nature's own landscaping has provided a border of alder, viburnum, sweet fern, and juniper with seasonally changing accents of bright flowers, or of fruits hanging in jeweled clusters in the fall. The road had no heavy load of traffic to support; there were few sharp curves or intersections where brush could obstruct the driver's vision. But the sprayers took over and the miles along that road became something to be traversed quickly, a sight to be endured with one's mind closed to thoughts of the sterile and hideous world we are letting our technicians make. But here and there authority had somehow faltered and by an unaccountable oversight there were oases of beauty in the midst of austere and regimented control---oases that made the desecration of the greater part of the road the more unbearable. In such places my spirit lifted to the sight of the drifts of white clover or the clouds of purple vetch with here and there the flaming cup of a wood lily.
\par
\noindent\textbf{中文。} 我很熟悉一段道路，在那里，自然自己的造景沿路铺设了一道桤木、荚蒾、香蕨木和刺柏的边界，季节变换时，明亮的花朵点缀其间；到了秋天，成串果实像珠宝般悬挂。那条路并不承载繁重交通；急弯和交叉路口也很少，灌丛几乎不会遮挡驾驶者视线。可是喷洒者接管之后，沿途数英里变成了人们只想尽快穿过的地方，成为一种必须忍受的景象；只有闭上心，不去想我们正让技术人员制造怎样一个贫瘠而丑陋的世界，才能承受。然而，这里那里，权威不知怎的失手了，由于某种无法解释的疏忽，在严苛而整齐划一的控制中留下了美的绿洲；这些绿洲反倒使道路大部分遭受的亵渎更加难以忍受。在这样的地方，看到一片片白三叶草，或一团团紫色野豌豆云雾，其间偶尔燃起一盏林百合的火焰杯盏，我的精神便随之振奋。
\par\medskip

\noindent\textbf{English.} Such plants are ``weeds'' only to those who make a business of selling and applying chemicals. In a volume of \textit{Proceedings} of one of the weed-control conferences that are now regular institutions, I once read an extraordinary statement of a weed killer's philosophy. The author defended the killing of good plants ``simply because they are in bad company.'' Those who complain about killing wildflowers along roadsides reminded him, he said, of antivivisectionists ``to whom, if one were to judge by their actions, the life of a stray dog is more sacred than the lives of children.''
\par
\noindent\textbf{中文。} 这些植物之所以是“杂草”，只是在那些以销售和施用化学品为业的人眼里如此。我曾在一本除草会议《论文集》中读到一段非同寻常的除草哲学宣言；这类会议如今已成为常设制度。作者为杀死好植物辩护，说理由“只不过是它们交了坏朋友”。他说，那些抱怨路边野花遭杀灭的人，让他想到反对动物活体实验者；“如果按他们的行为来判断，在他们眼中，一条流浪狗的生命比儿童的生命还神圣。”
\par\medskip

\noindent\textbf{English.} To the author of this paper, many of us would unquestionably be suspect, convicted of some deep perversion of character because we prefer the sight of the vetch and the clover and the wood lily in all their delicate and transient beauty to that of roadsides scorched as by fire, the shrubs brown and brittle, the bracken that once lifted high its proud lacework now withered and drooping. We would seem deplorably weak that we can tolerate the sight of such ``weeds,'' that we do not rejoice in their eradication, that we are not filled with exultation that man has once more triumphed over miscreant nature.
\par
\noindent\textbf{中文。} 在这篇论文的作者看来，我们许多人无疑都很可疑，必定会因某种深层性格败坏而被定罪：因为我们宁愿看见野豌豆、白三叶草和林百合那纤细而短暂的美，也不愿看见道路两旁像被火烧过一样，灌木褐黄而脆裂，曾经高举骄傲蕾丝花纹的蕨草如今枯萎垂落。我们似乎软弱得可悲，竟能容忍这些“杂草”的景象，竟不为它们被根除而欢欣，竟不会因人类又一次战胜作恶的自然而充满狂喜。
\par\medskip

\noindent\textbf{English.} Justice Douglas tells of attending a meeting of federal field men who were discussing protests by citizens against plans for the spraying of sagebrush that I mentioned earlier in this chapter. These men considered it hilariously funny that an old lady had opposed the plan because the wildflowers would be destroyed. ``Yet, was not her right to search out a banded cup or a tiger lily as inalienable as the right of stockmen to search out grass or of a lumberman to claim a tree?'' asks this humane and perceptive jurist. ``The esthetic values of the wilderness are as much our inheritance as the veins of copper and gold in our hills and the forests in our mountains.''
\par
\noindent\textbf{中文。} 道格拉斯大法官讲述过，他参加过一次联邦外勤人员会议，会上讨论公民对本章前面提到的喷洒蒿灌木计划的抗议。那些人觉得，一位老妇人因为野花会被毁灭而反对该计划，这件事滑稽极了。这位富有人道精神且洞察敏锐的法学家问道：“然而，她寻找斑纹杯花或虎百合的权利，难道不和牧场主寻找牧草的权利、木材商要求一棵树的权利同样不可剥夺吗？”“荒野的审美价值，正如山中铜金矿脉和山上森林一样，都是我们的遗产。”
\par\medskip

\noindent\textbf{English.} There is of course more to the wish to preserve our roadside vegetation than even such esthetic considerations. In the economy of nature the natural vegetation has its essential place. Hedgerows along country roads and bordering fields provide food, cover, and nesting areas for birds and homes for many small animals. Of some 70 species of shrubs and vines that are typical roadside species in the eastern states alone, about 65 are important to wildlife as food.
\par
\noindent\textbf{中文。} 当然，保存路边植被的愿望并不只出于这些审美考虑。在自然的经济体系中，天然植被有其不可或缺的位置。乡间道路和田地边缘的篱丛，为鸟类提供食物、隐蔽处和筑巢地，也为许多小动物提供家园。仅在东部各州，典型路边灌木和藤本约有七十种，其中约六十五种是野生动物的重要食物。
\par\medskip

\noindent\textbf{English.} Such vegetation is also the habitat of wild bees and other pollinating insects. Man is more dependent on these wild pollinators than he usually realizes. Even the farmer himself seldom understands the value of wild bees and often participates in the very measures that rob him of their services. Some agricultural crops and many wild plants are partly or wholly dependent on the services of the native pollinating insects. Several hundred species of wild bees take part in the pollination of cultivated crops---100 species visiting the flowers of alfalfa alone. Without insect pollination, most of the soil-holding and soil-enriching plants of uncultivated areas would die out, with far-reaching consequences to the ecology of the whole region. Many herbs, shrubs, and trees of forests and range depend on native insects for their reproduction; without these plants many wild animals and range stock would find little food. Now clean cultivation and the chemical destruction of hedgerows and weeds are eliminating the last sanctuaries of these pollinating insects and breaking the threads that bind life to life.
\par
\noindent\textbf{中文。} 这类植被也是野蜂和其他传粉昆虫的栖息地。人类对这些野生传粉者的依赖，远超通常所知。即使农民本人也很少理解野蜂的价值，并且常常参与那些恰恰剥夺其服务的措施。一些农作物和许多野生植物，部分或完全依赖本地传粉昆虫。数百种野蜂参与栽培作物的授粉；仅访问苜蓿花的就有一百种。没有昆虫授粉，未开垦地区多数固土、肥土植物都会消失，对整个区域的生态造成深远后果。森林和牧地中的许多草本、灌木和乔木，都依赖本地昆虫完成繁殖；没有这些植物，许多野生动物和牧场牲畜就几乎找不到食物。如今，清耕以及对篱丛和杂草的化学摧毁，正在消灭这些传粉昆虫最后的避难所，并切断生命彼此相系的丝线。
\par\medskip

\noindent\textbf{English.} These insects, so essential to our agriculture and indeed to our landscape as we know it, deserve something better from us than the senseless destruction of their habitat. Honeybees and wild bees depend heavily on such ``weeds'' as goldenrod, mustard, and dandelions for pollen that serves as the food of their young. Vetch furnishes essential spring forage for bees before the alfalfa is in bloom, tiding them over this early season so that they are ready to pollinate the alfalfa. In the fall they depend on goldenrod at a season when no other food is available, to stock up for the winter. By the precise and delicate timing that is nature's own, the emergence of one species of wild bees takes place on the very day of the opening of the willow blossoms. There is no dearth of men who understand these things, but these are not the men who order the wholesale drenching of the landscape with chemicals.
\par
\noindent\textbf{中文。} 这些昆虫对我们的农业，乃至对我们所熟知的整个景观都至关重要；我们给它们的，不应是对其栖息地的无谓毁灭。蜜蜂和野蜂高度依赖秋麒麟草、芥菜、蒲公英这类“杂草”的花粉，以此喂养幼虫。苜蓿开花前，野豌豆为蜂类提供必不可少的春季饲源，帮助它们度过早春，使其准备好为苜蓿授粉。秋天，在别无食物可得的季节，它们依赖秋麒麟草为越冬储备。按照自然自身精确而细腻的时间安排，某一种野蜂的羽化，正发生在柳花开放的那一天。理解这些事的人并不少；可是下令用化学品大规模浇灌景观的人，并不是他们。
\par\medskip

\noindent\textbf{English.} And where are the men who supposedly understand the value of proper habitat for the preservation of wildlife? Too many of them are to be found defending herbicides as ``harmless'' to wildlife because they are thought to be less toxic than insecticides. Therefore, it is said, no harm is done. But as the herbicides rain down on forest and field, on marsh and rangeland, they are bringing about marked changes and even permanent destruction of wildlife habitat. To destroy the homes and the food of wildlife is perhaps worse in the long run than direct killing.
\par
\noindent\textbf{中文。} 那些本应理解适当栖息地对保护野生动物有何价值的人，又在哪里？太多这样的人在为除草剂辩护，声称它们对野生动物“无害”，理由是除草剂被认为不如杀虫剂毒性强。因此，据说并未造成伤害。可是，当除草剂像雨一样落到森林和田野、沼泽和牧地上时，它们正在引起野生动物栖息地的显著改变，甚至造成永久性毁坏。摧毁野生动物的家园和食物，从长远看也许比直接杀戮更糟。
\par\medskip

\noindent\textbf{English.} The irony of this all-out chemical assault on roadsides and utility rights-of-way is twofold. It is perpetuating the problem it seeks to correct, for as experience has clearly shown, the blanket application of herbicides does not permanently control roadside ``brush'' and the spraying has to be repeated year after year. And as a further irony, we persist in doing this despite the fact that a perfectly sound method of selective spraying is known, which can achieve long-term vegetational control and eliminate repeated spraying in most types of vegetation.
\par
\noindent\textbf{中文。} 这种对路边和公用设施通道发动的全面化学攻击，具有双重讽刺。它延续了自己试图纠正的问题；经验已经清楚表明，除草剂的全面覆盖施用并不能永久控制路边“灌丛”，喷洒不得不年复一年重复。更讽刺的是，尽管我们明知一种完全稳妥的选择性喷洒方法存在，能够实现长期植被控制，并在大多数植被类型中免除反复喷洒，我们仍继续这样做。
\par\medskip

\noindent\textbf{English.} The object of brush control along roads and rights-of-way is not to sweep the land clear of everything but grass; it is, rather, to eliminate plants ultimately tall enough to present an obstruction to drivers' vision or interference with wires on rights-of-way. This means, in general, trees. Most shrubs are low enough to present no hazard; so, certainly, are ferns and wildflowers.
\par
\noindent\textbf{中文。} 道路和通道沿线灌丛控制的目标，并不是把土地扫荡到只剩青草；相反，它是要清除那些最终会长到足以遮挡驾驶者视线，或干扰通道上方线缆的植物。一般说来，这意味着树木。大多数灌木足够低，不构成危险；蕨类和野花当然更是如此。
\par\medskip

\noindent\textbf{English.} Selective spraying was developed by Dr. Frank Egler during a period of years at the American Museum of Natural History as director of a Committee for Brush Control Recommendations for Rights-of-Way. It took advantage of the inherent stability of nature, building on the fact that most communities of shrubs are strongly resistant to invasion by trees. By comparison, grasslands are easily invaded by tree seedlings. The object of selective spraying is not to produce grass on roadsides and rights-of-way but to eliminate the tall woody plants by direct treatment and to preserve all other vegetation. One treatment may be sufficient, with a possible follow-up for extremely resistant species; thereafter the shrubs assert control and the trees do not return. The best and cheapest controls for vegetation are not chemicals but other plants.
\par
\noindent\textbf{中文。} 选择性喷洒是弗兰克・埃格勒博士在美国自然历史博物馆任通道灌丛控制建议委员会主任期间，经多年发展出来的方法。它利用自然内在的稳定性，依据的事实是：大多数灌木生物群落强烈抵抗树木入侵。相比之下，草地很容易被树苗侵入。选择性喷洒的目标不是在路边和通道上生产草，而是通过直接处理清除高大的木本植物，并保存所有其他植被。一次处理也许就足够，对特别顽强的物种可作后续处理；此后灌木便取得控制，树木不会返回。控制植被最好、最便宜的办法不是化学品，而是其他植物。
\par\medskip

\noindent\textbf{English.} The method has been tested in research areas scattered throughout the eastern United States. Results show that once properly treated, an area becomes stabilized, requiring no re-spraying for at least 20 years. The spraying can often be done by men on foot, using knapsack sprayers, and having complete control over their material. Sometimes compressor pumps and material can be mounted on truck chassis, but there is no blanket spraying. Treatment is directed only to trees and any exceptionally tall shrubs that must be eliminated. The integrity of the environment is thereby preserved, the enormous value of the wildlife habitat remains intact, and the beauty of shrub and fern and wildflower has not been sacrificed.
\par
\noindent\textbf{中文。} 这种方法已在美国东部各地分散的研究区域中进行测试。结果表明，一片区域一旦经过适当处理，就会稳定下来，至少二十年不需要重新喷洒。喷洒常可由步行人员使用背负式喷雾器完成，从而完全控制药剂。有时也可把压缩泵和药剂装在卡车底盘上，但绝不进行覆盖式喷洒。处理只针对必须清除的树木和任何特别高大的灌木。环境完整性因而得以保存，野生动物栖息地的巨大价值保持无损，灌木、蕨类和野花之美也没有被牺牲。
\par\medskip

\noindent\textbf{English.} Here and there the method of vegetation management by selective spraying has been adopted. For the most part, entrenched custom dies hard and blanket spraying continues to thrive, to exact its heavy annual costs from the taxpayer, and to inflict its damage on the ecological web of life. It thrives, surely, only because the facts are not known. When taxpayers understand that the bill for spraying the town roads should come due only once a generation instead of once a year, they will surely rise up and demand a change of method.
\par
\noindent\textbf{中文。} 在一些地方，选择性喷洒这种植被管理方法已经被采用。但在大多数地方，根深蒂固的习惯难以消亡，全面喷洒仍然兴盛，年复一年向纳税人索取沉重成本，并把损害加诸生命的生态之网。它之所以兴盛，当然只是因为事实还不为人知。当纳税人明白，城镇道路喷洒的账单本应一代人才到期一次，而不是一年一次，他们必定会起来要求改变方法。
\par\medskip

\noindent\textbf{English.} Among the many advantages of selective spraying is the fact that it minimizes the amount of chemical applied to the landscape. There is no broadcasting of material but, rather, concentrated application to the base of the trees. The potential harm to wildlife is therefore kept to a minimum.
\par
\noindent\textbf{中文。} 选择性喷洒有许多优点，其中之一是它把施加到景观上的化学品数量降至最低。药剂不是被广泛撒布，而是集中施用于树基。因此，对野生动物的潜在伤害也被控制在最低限度。
\par\medskip

\noindent\textbf{English.} The most widely used herbicides are 2,4-D, 2,4,5-T, and related compounds. Whether or not these are actually toxic is a matter of controversy. People spraying their lawns with 2,4-D and becoming wet with spray have occasionally developed severe neuritis and even paralysis. Although such incidents are apparently uncommon, medical authorities advise caution in use of such compounds. Other hazards, more obscure, may also attend the use of 2,4-D. It has been shown experimentally to disturb the basic physiological process of respiration in the cell, and to imitate X-rays in damaging the chromosomes. Some very recent work indicates that reproduction of birds may be adversely affected by these and certain other herbicides at levels far below those that cause death.
\par
\noindent\textbf{中文。} 使用最广泛的除草剂是2,4-D、2,4,5-T以及相关化合物。它们是否真正有毒，仍有争议。有些人在用2,4-D喷洒自家草坪、并被喷雾淋湿后，偶尔发生严重神经炎，甚至瘫痪。尽管这类事件显然并不常见，医学权威仍建议谨慎使用这些化合物。使用2,4-D还可能伴随其他更隐蔽的危险。实验显示，它会干扰细胞呼吸这一基本生理过程，并会像X射线一样损伤染色体。最近的一些研究表明，在远低于致死剂量的水平上，这些和某些其他除草剂可能已经对鸟类繁殖造成不利影响。
\par\medskip

\noindent\textbf{English.} Apart from any directly toxic effects, curious indirect results follow the use of certain herbicides. It has been found that animals, both wild herbivores and livestock, are sometimes strangely attracted to a plant that has been sprayed, even though it is not one of their natural foods. If a highly poisonous herbicide such as arsenic has been used, this intense desire to reach the wilting vegetation inevitably has disastrous results. Fatal results may follow, also, from less toxic herbicides if the plant itself happens to be poisonous or perhaps to possess thorns or burs. Poisonous range weeds, for example, have suddenly become attractive to livestock after spraying, and the animals have died from indulging this unnatural appetite. The literature of veterinary medicine abounds in similar examples: swine eating sprayed cockleburs with consequent severe illness, lambs eating sprayed thistles, bees poisoned by pasturing on mustard sprayed after it came into bloom. Wild cherry, the leaves of which are highly poisonous, has exerted a fatal attraction for cattle once its foliage has been sprayed with 2,4-D. Apparently the wilting that follows spraying (or cutting) makes the plant attractive. Ragwort has provided other examples. Livestock ordinarily avoid this plant unless forced to turn to it in late winter and early spring by lack of other forage. However, the animals eagerly feed on it after its foliage has been sprayed with 2,4-D.
\par
\noindent\textbf{中文。} 除了任何直接毒性影响之外，某些除草剂的使用还会带来奇特的间接结果。人们发现，无论野生食草动物还是家畜，有时都会奇怪地被喷洒过的植物吸引，哪怕那并不是它们的天然食物。如果使用的是砷这类剧毒除草剂，这种接近萎蔫植被的强烈欲望，必然造成灾难性后果。即使除草剂毒性较低，只要植物本身恰好有毒，或者也许带刺或带芒刺，也可能导致死亡。例如，有毒的牧场杂草在喷洒后会突然变得吸引牲畜，动物因纵容这种反常食欲而死去。兽医学文献中充满类似例子：猪吃了喷洒过的苍耳而重病；羔羊吃了喷洒过的蓟；蜜蜂在芥菜开花后取食被喷洒的芥菜而中毒。野樱桃的叶子毒性很强；一旦其叶片被2,4-D喷洒，就会对牛产生致命吸引。显然，喷洒或切割之后发生的萎蔫，使植物变得有吸引力。千里光也提供了其他例子。牲畜通常避开这种植物，除非在晚冬和早春因缺乏其他饲草而被迫转向它。然而，它的叶片一经2,4-D喷洒，动物便会急切取食。
\par\medskip

\noindent\textbf{English.} The explanation of this peculiar behavior sometimes appears to lie in the changes which the chemical brings about in the metabolism of the plant itself. There is temporarily a marked increase in sugar content, making the plant more attractive to many animals.
\par
\noindent\textbf{中文。} 对这种奇特行为的解释，有时似乎在于化学品对植物自身代谢造成的变化。植物含糖量会暂时显著增加，从而对许多动物更具吸引力。
\par\medskip

\noindent\textbf{English.} Another curious effect of 2,4-D has important effects for livestock, wildlife, and apparently for men as well. Experiments carried out about a decade ago showed that after treatment with this chemical there is a sharp increase in the nitrate content of corn and of sugar beets. The same effect was suspected in sorghum, sunflower, spiderwort, lamb's quarters, pigweed, and smartweed. Some of these are normally ignored by cattle, but are eaten with relish after treatment with 2,4-D. A number of deaths among cattle have been traced to sprayed weeds, according to some agricultural specialists. The danger lies in the increase in nitrates, for the peculiar physiology of the ruminant at once poses a critical problem. Most such animals have a digestive system of extraordinary complexity, including a stomach divided into four chambers. The digestion of cellulose is accomplished through the action of microorganisms (rumen bacteria) in one of the chambers. When the animal feeds on vegetation containing an abnormally high level of nitrates, the microorganisms in the rumen act on the nitrates to change them into highly toxic nitrites. Thereafter a fatal chain of events ensues: the nitrites act on the blood pigment to form a chocolate-brown substance in which the oxygen is so firmly held that it cannot take part in respiration, hence oxygen is not transferred from the lungs to the tissues. Death occurs within a few hours from anoxia, or lack of oxygen. The various reports of livestock losses after grazing on certain weeds treated with 2,4-D therefore have a logical explanation. The same danger exists for wild animals belonging to the group of ruminants, such as deer, antelope, sheep, and goats.
\par
\noindent\textbf{中文。} 2,4-D的另一种奇特影响，对牲畜、野生动物，显然也对人具有重要后果。约十年前进行的实验显示，经这种化学品处理后，玉米和甜菜中的硝酸盐含量急剧上升。同样的影响也被怀疑出现在高粱、向日葵、紫露草、藜、苋和蓼等植物中。其中一些植物通常被牛忽视，但经2,4-D处理后却会被津津有味地吃掉。一些农业专家认为，已有若干牛死亡事件可追溯到喷洒过的杂草。危险在于硝酸盐增加，因为反刍动物独特的生理立即带来关键问题。多数这类动物拥有极其复杂的消化系统，包括分成四个室的胃。纤维素消化是在其中一个胃室里通过微生物（瘤胃细菌）的作用完成的。当动物取食含有异常高水平硝酸盐的植被时，瘤胃中的微生物作用于硝酸盐，将其转化为剧毒的亚硝酸盐。随后，一连串致命事件发生：亚硝酸盐作用于血液色素，形成一种巧克力褐色物质，氧被牢牢结合其中，以致不能参与呼吸；因此，氧不能从肺转移到组织。数小时内，动物便因缺氧而死亡。因此，关于牲畜在取食经2,4-D处理的某些杂草后死亡的各种报告，有了合乎逻辑的解释。属于反刍类的野生动物，如鹿、叉角羚、绵羊和山羊，也面临同样危险。
\par\medskip

\noindent\textbf{English.} Although various factors (such as exceptionally dry weather) can cause an increase in nitrate content, the effect of the soaring sales and applications of 2,4-D cannot be ignored. The situation was considered important enough by the University of Wisconsin Agricultural Experiment Station to justify a warning in 1957 that ``plants killed by 2,4-D may contain large amounts of nitrate.'' The hazard extends to human beings as well as animals and may help to explain the recent mysterious increase in ``silo deaths.'' When corn, oats, or sorghum containing large amounts of nitrates are ensiled they release poisonous nitrogen oxide gases, creating a deadly hazard to anyone entering the silo. Only a few breaths of one of these gases can cause a diffuse chemical pneumonia. In a series of such cases studied by the University of Minnesota Medical School all but one terminated fatally.
\par
\noindent\textbf{中文。} 虽然多种因素（例如异常干旱天气）都能导致硝酸盐含量上升，但2,4-D销量和施用量飙升的影响不能被忽视。威斯康星大学农业试验站认为情况足够重要，因而在1957年发出警告：“被2,4-D杀死的植物可能含有大量硝酸盐。”这种危险不仅延伸到动物，也延伸到人类，并可能有助于解释近来“青贮窖死亡”事件神秘增加的原因。当含有大量硝酸盐的玉米、燕麦或高粱被青贮时，会释放有毒的氮氧化物气体，对任何进入青贮窖的人构成致命危险。只要吸入几口其中一种气体，就可能造成弥漫性化学性肺炎。明尼苏达大学医学院研究的一系列此类病例中，除一例外全部致死。
\par\medskip

\noindent\textbf{English.} ``Once again we are walking in nature like an elephant in the china cabinet.'' So C. J. Briejer, a Dutch scientist of rare understanding, sums up our use of weed killers. ``In my opinion too much is taken for granted. We do not know whether all weeds in crops are harmful or whether some of them are useful,'' says Dr. Briejer.
\par
\noindent\textbf{中文。} “我们又一次像大象走进瓷器柜那样行走在自然中。”荷兰科学家C・J・布里耶尔以罕见的理解力，这样总结我们对除草剂的使用。布里耶尔博士说：“在我看来，人们把太多事情视为理所当然。我们并不知道作物中的所有杂草是否都有害，也不知道其中有些是否有益。”
\par\medskip

\noindent\textbf{English.} Seldom is the question asked: What is the relation between the weed and the soil? Perhaps, even from our narrow standpoint of direct self-interest, the relation is a useful one. As we have seen, soil and the living things in and upon it exist in a relation of interdependence and mutual benefit. Presumably the weed is taking something from the soil; perhaps it is also contributing something to it. A practical example was provided recently by the parks in a city in Holland. The roses were doing badly. Soil samples showed heavy infestations by tiny nematode worms. Scientists of the Dutch Plant Protection Service did not recommend chemical sprays or soil treatments; instead, they suggested that marigolds be planted among the roses. This plant, which the purist would doubtless consider a weed in any rosebed, releases an excretion from its roots that kills the soil nematodes. The advice was taken; some beds were planted with marigolds, some left without as controls. The results were striking. With the aid of the marigolds the roses flourished; in the control beds they were sickly and drooping. Marigolds are now used in many places for combating nematodes.
\par
\noindent\textbf{中文。} 人们很少提出这样的问题：杂草与土壤之间有什么关系？也许，即使从我们狭隘的直接自利角度看，这种关系也是有用的。正如我们已经看到的，土壤以及生活在其中、其上的生命，存在相互依赖、互利共生的关系。可以推想，杂草从土壤中取走某些东西；也许它也向土壤贡献某些东西。荷兰一座城市的公园最近提供了一个实际例子。玫瑰长势很差。土样显示微小线虫严重侵染。荷兰植物保护局的科学家没有建议化学喷洒或土壤处理；相反，他们建议在玫瑰之间种植万寿菊。这种植物，洁癖式的园艺家无疑会认为它在任何玫瑰花坛里都是杂草，却会从根部分泌一种物质，杀死土壤线虫。建议被采纳了；一些花坛种上万寿菊，另一些不种，作为对照。结果非常显著。在万寿菊帮助下，玫瑰繁茂生长；对照花坛中的玫瑰则病弱下垂。如今，万寿菊已在许多地方用于防治线虫。
\par\medskip

\noindent\textbf{English.} In the same way, and perhaps quite unknown to us, other plants that we ruthlessly eradicate may be performing a function that is necessary to the health of the soil. One very useful function of natural plant communities---now pretty generally stigmatized as ``weeds''---is to serve as an indicator of the condition of the soil. This useful function is of course lost where chemical weed killers have been used.
\par
\noindent\textbf{中文。} 同样地，也许我们全然不知，那些被我们无情根除的其他植物，可能正在履行一种维持土壤健康所必需的功能。天然植物生物群落如今已普遍被污名化为“杂草”，但它们有一项非常有用的功能：指示土壤状况。当然，在使用化学除草剂的地方，这种有用功能便丧失了。
\par\medskip

\noindent\textbf{English.} Those who find an answer to all problems in spraying also overlook a matter of great scientific importance---the need to preserve some natural plant communities. We need these as a standard against which we can measure the changes our own activities bring about. We need them as wild habitats in which original populations of insects and other organisms can be maintained, for, as will be explained in Chapter 16, the development of resistance to insecticides is changing the genetic factors of insects and perhaps other organisms. One scientist has even suggested that some sort of ``zoo'' should be established to preserve insects, mites, and the like, before their genetic composition is further changed.
\par
\noindent\textbf{中文。} 那些认为喷洒能回答一切问题的人，也忽视了一件具有重大科学意义的事：需要保存某些天然植物生物群落。我们需要它们作为标准，用来衡量我们自身活动带来的变化。我们还需要它们作为野生栖息地，使昆虫和其他生物的原初种群得以维持；因为正如第十六章将说明的，对杀虫剂抗性的产生，正在改变昆虫，也许还有其他生物的遗传因素。一位科学家甚至建议，在昆虫、螨类等的遗传组成进一步改变之前，应建立某种“动物园”来保存它们。
\par\medskip

\noindent\textbf{English.} Some experts warn of subtle but far-reaching vegetational shifts as a result of the growing use of herbicides. The chemical 2,4-D, by killing out the broad-leaved plants, allows the grasses to thrive in the reduced competition---now some of the grasses themselves have become ``weeds,'' presenting a new problem in control and giving the cycle another turn. This strange situation is acknowledged in a recent issue of a journal devoted to crop problems: ``With the widespread use of 2,4-D to control broadleaved weeds, grass weeds in particular have increasingly become a threat to corn and soybean yields.''
\par
\noindent\textbf{中文。} 一些专家警告，除草剂使用日增，可能导致微妙却影响深远的植被转变。2,4-D这种化学品通过杀灭阔叶植物，减少竞争，使禾草繁盛；如今，一些禾草本身又变成了“杂草”，提出新的控制问题，使循环又转了一圈。最近一期专门讨论作物问题的期刊承认了这种奇怪状况：“随着2,4-D被广泛用于控制阔叶杂草，禾本科杂草尤其日益成为对玉米和大豆产量的威胁。”
\par\medskip

\noindent\textbf{English.} Ragweed, the bane of hay fever sufferers, offers an interesting example of the way efforts to control nature sometimes boomerang. Many thousands of gallons of chemicals have been discharged along roadsides in the name of ragweed control. But the unfortunate truth is that blanket spraying is resulting in more ragweed, not less. Ragweed is an annual; its seedlings require open soil to become established each year. Our best protection against this plant is therefore the maintenance of dense shrubs, ferns, and other perennial vegetation. Spraying frequently destroys this protective vegetation and creates open, barren areas which the ragweed hastens to fill. It is probable, moreover, that the pollen content of the atmosphere is not related to roadside ragweed, but to the ragweed of city lots and fallow fields.
\par
\noindent\textbf{中文。} 豚草是花粉热患者的祸害，它提供了一个有趣例子，说明控制自然的努力有时会反噬。以控制豚草为名，成千上万加仑化学品被排放到道路两旁。但不幸的事实是，覆盖式喷洒带来的豚草不是更少，而是更多。豚草是一年生植物；它的幼苗每年都需要裸露土壤才能定居。因此，我们抵御这种植物的最好保护，是维持浓密灌木、蕨类和其他多年生植被。喷洒往往破坏这种保护性植被，制造开阔贫瘠的区域，而豚草会迅速填补。况且，大气中的花粉含量很可能并不与路边豚草有关，而是与城市空地和休耕地上的豚草有关。
\par\medskip

\noindent\textbf{English.} The booming sales of chemical crabgrass killers are another example of how readily unsound methods catch on. There is a cheaper and better way to remove crabgrass than to attempt year after year to kill it out with chemicals. This is to give it competition of a kind it cannot survive, the competition of other grass. Crabgrass exists only in an unhealthy lawn. It is a symptom, not a disease in itself. By providing a fertile soil and giving the desired grasses a good start, it is possible to create an environment in which crabgrass cannot grow, for it requires open space in which it can start from seed year after year.
\par
\noindent\textbf{中文。} 化学马唐除草剂销量激增，是不健全方法多么容易流行的另一个例子。要清除马唐，有一种比年复一年试图用化学品杀死它更便宜、更好的办法：给予它无法承受的竞争，也就是其他草类的竞争。马唐只存在于不健康的草坪中。它是一种症状，本身并不是疾病。只要提供肥沃土壤，并让所需草类有良好开端，就可以创造一种马唐不能生长的环境，因为它需要开阔空间，才能年复一年从种子开始生长。
\par\medskip

\noindent\textbf{English.} Instead of treating the basic condition, suburbanites---advised by nurserymen who in turn have been advised by the chemical manufacturers---continue to apply truly astonishing amounts of crabgrass killers to their lawns each year. Marketed under trade names which give no hint of their nature, many of these preparations contain such poisons as mercury, arsenic, and chlordane. Application at the recommended rates leaves tremendous amounts of these chemicals on the lawn. Users of one product, for example, apply 60 pounds of technical chlordane to the acre if they follow directions. If they use another of the many available products, they are applying 175 pounds of metallic arsenic to the acre. The toll of dead birds, as we shall see in Chapter 8, is distressing. How lethal these lawns may be for human beings is unknown.
\par
\noindent\textbf{中文。} 郊区居民没有处理根本状况，而是在苗圃业者建议下继续每年向草坪施用数量惊人的马唐除草剂；而这些苗圃业者又接受了化学品制造商的建议。许多这类制剂以商标名称销售，名称丝毫不透露其本质，却含有汞、砷、氯丹等毒物。按推荐剂量施用，会在草坪上留下大量化学品。例如，某一产品的用户若按说明操作，每英亩便施用了六十磅工业氯丹。若使用另一些可选产品中的一种，他们每英亩施用的则是一百七十五磅金属砷。正如我们将在第八章看到的，死亡鸟类的代价令人痛心。这些草坪对人类可能有多致命，目前尚无人知晓。
\par\medskip

\noindent\textbf{English.} The success of selective spraying for roadside and right-of-way vegetation, where it has been practiced, offers hope that equally sound ecological methods may be developed for other vegetation programs for farms, forests, and ranges---methods aimed not at destroying a particular species but at managing vegetation as a living community.
\par
\noindent\textbf{中文。} 在已经实践的地方，路边和通道植被选择性喷洒的成功，为农场、森林和牧地的其他植被计划提供了希望：同样健全的生态方法也许能够发展出来。这些方法的目标不是毁灭某个特定物种，而是把植被作为一个有生命的生物群落来管理。
\par\medskip

\noindent\textbf{English.} Other solid achievements show what can be done. Biological control has achieved some of its most spectacular successes in the area of curbing unwanted vegetation. Nature herself has met many of the problems that now beset us, and she has usually solved them in her own successful way. Where man has been intelligent enough to observe and to emulate Nature he, too, is often rewarded with success.
\par
\noindent\textbf{中文。} 其他扎实成就显示了可以做到什么。生物防治在遏制不受欢迎植被的领域，取得了一些最引人注目的成功。自然本身已经遭遇过许多如今困扰我们的问题，并且通常以自己的成功方式解决了它们。当人类足够聪明，去观察并效法自然时，他也常常得到成功的回报。
\par\medskip

\noindent\textbf{English.} An outstanding example in the field of controlling unwanted plants is the handling of the Klamath-weed problem in California. Although the Klamath weed, or goatweed, is a native of Europe (where it is called St. Johnswort), it accompanied man in his westward migrations, first appearing in the United States in 1793 near Lancaster, Pennsylvania. By 1900 it had reached California in the vicinity of the Klamath River, hence the name locally given to it. By 1929 it had occupied about 100,000 acres of rangeland, and by 1952 it had invaded some two and one half million acres.
\par
\noindent\textbf{中文。} 在控制不受欢迎植物的领域，加利福尼亚处理克拉马斯草问题，是一个杰出例子。克拉马斯草，又称山羊草，原产欧洲（在那里称为圣约翰草或贯叶连翘），却伴随人类西迁，1793年首先出现在美国宾夕法尼亚州兰开斯特附近。到1900年，它已抵达加利福尼亚克拉马斯河一带，因此获得当地名称。到1929年，它占据了约十万英亩牧地；到1952年，入侵面积已达约二百五十万英亩。
\par\medskip

\noindent\textbf{English.} Klamath weed, quite unlike such native plants as sagebrush, has no place in the ecology of the region, and no animals or other plants require its presence. On the contrary, wherever it appeared livestock became ``scabby, sore-mouthed, and unthrifty'' from feeding on this toxic plant. Land values declined accordingly, for the Klamath weed was considered to hold the first mortgage.
\par
\noindent\textbf{中文。} 克拉马斯草与蒿灌木这类本地植物截然不同，在该地区生态中没有位置，也没有动物或其他植物需要它存在。相反，它出现在哪里，牲畜因取食这种有毒植物就会变得“结痂、口腔溃痛、发育不良”。土地价值也相应下降，因为克拉马斯草被认为握有第一抵押权。
\par\medskip

\noindent\textbf{English.} In Europe the Klamath weed, or St. Johnswort, has never become a problem because along with the plant there have developed various species of insects; these feed on it so extensively that its abundance is severely limited. In particular, two species of beetles in southern France, pea-sized and of metallic color, have their whole beings so adapted to the presence of the weed that they feed and reproduce only upon it.
\par
\noindent\textbf{中文。} 在欧洲，克拉马斯草，也就是圣约翰草，从未成为问题，因为与这种植物一起演化出了多种昆虫；这些昆虫大量取食它，使其丰度受到严格限制。尤其是在法国南部，有两种豌豆大小、带金属色泽的甲虫，其整个生命都如此适应这种杂草的存在，以至于只在它身上取食和繁殖。
\par\medskip

\noindent\textbf{English.} It was an event of historic importance when the first shipments of these beetles were brought to the United States in 1944, for this was the first attempt in North America to control a plant with a plant-eating insect. By 1948 both species had become so well established that no further importations were needed. Their spread was accomplished by collecting beetles from the original colonies and redistributing them at the rate of millions a year. Within small areas the beetles accomplish their own dispersion, moving on as soon as the Klamath weed dies out and locating new stands with great precision. And as the beetles thin out the weed, desirable range plants that have been crowded out are able to return.
\par
\noindent\textbf{中文。} 1944年，第一批这类甲虫被运到美国，这是一个具有历史意义的事件，因为这是北美首次尝试用食植昆虫控制一种植物。到1948年，两种甲虫都已建立稳固种群，不再需要进一步进口。它们的扩散，是通过从最初群落中采集甲虫，并以每年数百万只的速度重新分布来完成的。在小范围内，甲虫会自行扩散；克拉马斯草一旦死去，它们便迁往别处，并以极高精度找到新的草丛。随着甲虫削弱这种杂草，那些被挤出的有益牧地植物便能够返回。
\par\medskip

\noindent\textbf{English.} A ten-year survey completed in 1959 showed that control of the Klamath weed had been ``more effective than hoped for even by enthusiasts,'' with the weed reduced to a mere 1 per cent of its former abundance. This token infestation is harmless and is actually needed in order to maintain a population of beetles as protection against a future increase in the weed.
\par
\noindent\textbf{中文。} 1959年完成的一项十年调查显示，对克拉马斯草的控制“比热心者所希望的还要有效”，这种杂草已减少到原来丰度的仅百分之一。这种象征性的侵染并无害处，实际上还为维持甲虫种群所必需，以防该杂草未来再次增加。
\par\medskip

\noindent\textbf{English.} Another extraordinarily successful and economical example of weed control may be found in Australia. With the colonists' usual taste for carrying plants or animals into a new country, a Captain Arthur Phillip had brought various species of cactus into Australia about 1787, intending to use them in culturing cochineal insects for dye. Some of the cacti or prickly pears escaped from his gardens and by 1925 about 20 species could be found growing wild. Having no natural controls in this new territory, they spread prodigiously, eventually occupying about 60 million acres. At least half of this land was so densely covered as to be useless.
\par
\noindent\textbf{中文。} 另一个极其成功且经济的杂草控制例子，可在澳大利亚找到。殖民者惯常喜欢把植物或动物带入新国土；约在1787年，亚瑟・菲利普船长把多种仙人掌带到澳大利亚，打算用它们培养胭脂虫制取染料。其中一些仙人掌或梨果仙人掌从他的花园逃逸；到1925年，约二十个种已在野外生长。在这片新领地缺乏自然控制因素，它们便异常扩散，最终占据约六千万英亩土地。其中至少一半覆盖过密，已无法利用。
\par\medskip

\noindent\textbf{English.} In 1920 Australian entomologists were sent to North and South America to study insect enemies of the prickly pears in their native habitat. After trials of several species, 3 billion eggs of an Argentine moth were released in Australia in 1930. Seven years later the last dense growth of the prickly pear had been destroyed and the once uninhabitable areas reopened to settlement and grazing. The whole operation had cost less than a penny per acre. In contrast, the unsatisfactory attempts at chemical control in earlier years had cost about \pounds 10 per acre.
\par
\noindent\textbf{中文。} 1920年，澳大利亚昆虫学家被派往南、北美洲，研究梨果仙人掌在原生栖息地中的昆虫天敌。试验若干物种之后，1930年，三十亿枚阿根廷蛾卵在澳大利亚释放。七年后，最后一片茂密的梨果仙人掌丛被摧毁，曾经不宜居住的地区重新向定居和放牧开放。整个行动每英亩花费不到一便士。相比之下，早年不令人满意的化学控制尝试，每英亩花费约十英镑。
\par\medskip

\noindent\textbf{English.} Both of these examples suggest that extremely effective control of many kinds of unwanted vegetation might be achieved by paying more attention to the role of plant-eating insects. The science of range management has largely ignored this possibility, although these insects are perhaps the most selective of all grazers and their highly restricted diets could easily be turned to man's advantage.
\par
\noindent\textbf{中文。} 这两个例子都表明，只要更加重视食植昆虫的作用，许多种不受欢迎植被或许都能得到极其有效的控制。牧地管理科学在很大程度上忽视了这种可能性，尽管这些昆虫也许是所有取食者中选择性最高的，其高度受限的食谱很容易转化为有利于人类的工具。
\par
